Fix \(P\in\mathcal P_{\mathrm{NG},n}\).  By the definitions of \(D_n\) and \(Z_n\),
\[
 Z_n=\eta+D_n,
\]
where \(D_n\) is measurable with respect to \(X\).  The independence restriction \(\eta\perp X\) therefore gives, for every \(z\in\mathbb C\),
\[
 F_n(z)
 =E_P[e^{z(\eta+D_n)}]
 =E_P[e^{z\eta}]E_P[e^{zD_n}]
 =M(z)H_n(z).
\]
All expectations are finite: \(D_n\) is bounded by \(2C_g\), while sub-Gaussianity of \(\eta\) gives finite exponential moments of every real order.  The same bounds justify differentiation under the expectation, so
\[
 F_n'(z)=E_P[Z_ne^{zZ_n}].
\]
The residual identity in the symbol definitions is
\[
 Y=\theta_0Z_n+b_n(X)+\xi.
\]
Because \(Z_n=T-\bar g_n(X)\) is measurable with respect to \((X,T)\), outcome mean independence yields
\[
 E_P[\xi e^{zZ_n}]
 =E_P\!\left[e^{zZ_n}E_P[\xi\mid X,T]\right]
 =0.
\]
Independence of \(\eta\) and \(X\) also gives
\[
 E_P[b_n(X)e^{zZ_n}]
 =E_P[e^{z\eta}]E_P[b_n(X)e^{zD_n}]
 =M(z)B_n(z).
\]
Combining the last three displays proves
\[
 G_n(z)=E_P[Ye^{zZ_n}]=\theta_0F_n'(z)+M(z)B_n(z).
\]
The domination already noted also shows that all four transforms in the identities are entire.